%! Author = sriadityavedantam
%! Date = 3/30/26

\section{Ideals and Quotient Rings}

\begin{definition}[Quotient Rings]
Congruence classes mod $f(x)$ are denoted by $\f[x]/(f(x))$, and this is a ring.
The additive identity of this ring is $[0]=[f(x)]$.
This together is called a quotient ring, and it is closed under addition and multiplication.
\end{definition}

\begin{theorem}[Class of $g(x)$]{
$[g(x)]\in \f[x]/(f(x))$ is a unit if and only if
\[
\gcd(f(x),g(x))=1.
\]
That is, $g(x)$ and $f(x)$ are relatively prime.
}
($\impliedby$).
Suppose $\gcd(f(x),g(x))=1$.
Then there exist $w(x),v(x)\in \f[x]$ such that
\[
w(x)g(x)+v(x)f(x)=1.
\]
Thus
\[
[w(x)g(x)]=[1],
\]
so
\[
[w(x)]=[g(x)]^{-1}.
\]

($\implies$).
Suppose $[g(x)]$ is a unit.
Then there exists $w(x)\in \f[x]$ such that
\[
[w(x)g(x)]=[1].
\]
So
\[
w(x)g(x)\equiv 1 \mod f(x),
\]
which means
\[
f(x)\mid w(x)g(x)-1.
\]
Therefore there exists $v(x)\in \f[x]$ such that
\begin{align*}
    w(x)g(x)-1 &= v(x)f(x),\\
    w(x)g(x)-v(x)f(x) &= 1.
\end{align*}
Therefore,
\[
\gcd(f(x),g(x))=1.
\]
\end{theorem}

\begin{corollary}{
If $f(x)$ is irreducible in $\f[x]$, then $\f[x]/(f(x))$ is a field.
}
If $g(x)\in \f[x]$ and $[g(x)]\neq [0]$, then $f(x)\nmid g(x)$.
Since $f(x)$ is irreducible, this implies
\[
\gcd(f(x),g(x))=1.
\]
So by the previous theorem, $[g(x)]$ is a unit in $\f[x]/(f(x))$.
Thus every nonzero element is a unit, and $\f[x]/(f(x))$ is a field.
\end{corollary}

Let
\[
\E=\fxmodfx.
\]
Then we have an injection from $\f\to\E$ where $a\mapsto [a]$.
So we may consider $\f$ as a subfield of $\E$.

\begin{definition}[Roots in Quotient Rings]
Suppose $\f\subseteq \E$.
Let $\alpha=[x]$.
Then for $f(x)\in \E[x]$, we have
\[
f(\alpha)=[0].
\]
\end{definition}

\begin{axiom}{
$\f\cong \e$.
}
\end{axiom}

\begin{definition}[Ideal]
Let $\ri$ be a commutative ring.
Given that $I\subseteq \ri$, we call $I$ an ideal if and only if $I$ is a subring of $\ri$ and whenever $r\in \ri$ and $a\in I$, then
\[
ra\in I.
\]
\end{definition}

\begin{definition}[Congruence mod $I$]
Suppose $r,s\in \ri$.
We write
\[
r\equiv s\mod I
\]
if
\[
r-s\in I.
\]
\end{definition}

\begin{theorem}[Congruence mod $I$ is an Equivalence Relation]{
Given $a,b,c\in \ri$, congruence mod $I$ is reflexive, symmetric, and transitive.
}
\textbf{Reflexive.}
\[
a\equiv a\mod I
\]
because
\[
a-a=0\in I.
\]

\textbf{Symmetric.}
If
\[
a\equiv b\mod I,
\]
then
\[
a-b\in I.
\]
Since $I$ is a subring, $-(a-b)=b-a\in I$, so
\[
b\equiv a\mod I.
\]

\textbf{Transitive.}
If
\[
a\equiv b\mod I
\quad\text{and}\quad
b\equiv c\mod I,
\]
then
\[
a-b\in I
\quad\text{and}\quad
b-c\in I.
\]
Since $I$ is closed under addition,
\[
(a-b)+(b-c)=a-c\in I.
\]
Thus
\[
a\equiv c\mod I.
\]
\end{theorem}

\begin{definition}[Coset]
Instead of $[a]$ for $a\mod I$, we have the notation
\[
a+I\coloneqq \{a+i:i\in I\},
\]
called a coset.
\end{definition}

For example,
\[
[a]_{\z_m}=a+m\z.
\]

\begin{definition}[Quotient Ring]
$\ri/I$ is called a quotient ring.
\end{definition}

\begin{theorem}[Addition on Ideals]{
\[
(a+I)+(b+I)=a+b+I
\]
and
\[
(a+I)(b+I)=ab+I.
\]
}
Suppose $a+I=c+I$ and $b+I=d+I$.
Then
\[
c-a\in I
\quad\text{and}\quad
d-b\in I.
\]
Hence
\[
(c-a)+(d-b)\in I,
\]
which implies
\[
(c+d)-(a+b)\in I.
\]
Therefore
\[
c+d+I=a+b+I.
\]
To prove multiplication, since $c-a,d-b\in I$, we have
\[
c(d-b)\in I
\quad\text{and}\quad
b(c-a)\in I
\]
by the absorption property of ideals.
Thus
\[
c(d-b)+b(c-a)\in I.
\]
But
\[
c(d-b)+b(c-a)=cd-cb+bc-ba=cd-ab.
\]
So
\[
cd-ab\in I,
\]
hence
\[
ab+I=cd+I.
\]
\end{theorem}

Quotient rings are independently associated with homomorphisms $\phi:\ri\to\si$.

\begin{definition}[Generators]
If $\ri$ is any commutative ring and $a\in \ri$, the ideal generated by $a$ is
\[
\{ra:r\in \ri\}\eqqcolon (a).
\]
\end{definition}

\begin{lemma}{
$(a)$ is an ideal of $\ri$.
}
\textbf{Case 1.}
If $r_{1}a,r_{2}a\in (a)$, then
\[
r_{1}a+r_{2}a=(r_1+r_2)a\in (a).
\]

\textbf{Case 2.}
If $ra\in (a)$ and $s\in \ri$, then
\[
s(ra)=(sr)a\in (a).
\]
\end{lemma}

These generators are called the principal ideal generated by $a$.

\begin{theorem}{
If $p(x)$ is irreducible in $\f[x]$, then the following are equivalent:
\[
\f[x]/(p(x)) \text{ is a field}
\iff
\f[x]/(p(x)) \text{ is an integral domain}.
\]
}
We already know that if $p(x)$ is irreducible, then $\f[x]/(p(x))$ is a field.
Every field is an integral domain, so one direction is immediate.
Conversely, if $\f[x]/(p(x))$ is an integral domain, then $(p(x))$ is a prime ideal.
In $\f[x]$, prime ideals generated by nonzero elements correspond to irreducible polynomials.
Hence $p(x)$ is irreducible.
\end{theorem}

Let $\ri$ be a commutative ring with $1\in \ri$.
Let $A$ be any subset of $\ri$.
The ideal generated by $A$ is the set of all finite linear combinations of elements of $A$:
\[
(A)\coloneqq \{r_{1}a_1+\hdots+r_{n}a_n:r_i\in \ri,\ a_i\in A\}.
\]
Then $(A)$ is the intersection of all ideals containing $A$.
Suppose $\ri=\z$ and $a,b\in \z$.
Then the ideal generated by $(a,b)$ is
\[
(a,b)\coloneqq \{xa+by:x,y\in \z\}.
\]
In fact,
\[
(a,b)=\{r\cdot \gcd(a,b):r\in \z\}.
\]
$\z$ and $\f[x]$ are called principal ideal domains, while $\z[x]$ and $\Q[x,y]$ are not principal ideal domains.
Let
\[
\phi:\z\to \z/10\z,
\]
defined by
\[
\phi(a)=[a]_{10}=a+10\z.
\]

\begin{definition}[Kernel]
Let
\[
K\coloneqq \{x\in \z:\phi(x)=0\}.
\]
This is called the kernel of $\phi$, written $\ker\phi$.
\end{definition}

\begin{theorem}{
$K$ is an ideal in $\ri$.
}
\textbf{(1)}
Suppose $x,y\in \ker\phi$.
Then
\[
\phi(x)=\phi(y)=0.
\]
So
\[
\phi(x+y)=\phi(x)+\phi(y)=0,
\]
hence $x+y\in \ker\phi$.

\textbf{(2)}
Suppose $x\in \ker\phi$ and $r\in \ri$.
Then
\[
\phi(rx)=\phi(r)\phi(x)=\phi(r)\cdot 0=0.
\]
So $rx\in \ker\phi$.
Therefore $\ker\phi$ is an ideal in $\ri$.
\end{theorem}

From the previous example,
\[
\ker\phi=(10)=10\z.
\]
What we learned prior is that
\[
\z/10\z \cong \z_{10}.
\]

\begin{definition}[Image]
\[
\Ima\phi \coloneqq \{s\in \si:\exists r\in \ri,\ \phi(r)=s\}.
\]
\end{definition}

\begin{theorem*}[First Isomorphism Theorem]{
Suppose $\phi:\ri\to\si$ is a homomorphism.
Let $K=\ker\phi$.
Define
\[
\overline{\phi}:\ri/K\to \Ima\phi
\]
by
\[
\overline{\phi}(r+K)=\phi(r).
\]
Then $\overline{\phi}$ is an isomorphism from $\ri/K$ to $\Ima\phi$.
So
\[
\ri/K\cong \Ima\phi.
\]
}
\end{theorem*}

\begin{proposition}{
Suppose $\phi:\ri\to\si$ is a ring homomorphism.
Then $\phi$ is injective if and only if
\[
\ker\phi=\{0\}.
\]
}
($\implies$).
Suppose $\phi$ is injective.
Let $r\in \ker\phi$.
So
\[
\phi(r)=0.
\]
But also
\[
\phi(0)=0.
\]
Since $\phi$ is injective, it follows that
\[
r=0.
\]
Thus
\[
\ker\phi=\{0\}.
\]

($\impliedby$).
Suppose
\[
\ker\phi=\{0\}.
\]
Let $r,s\in \ri$ with
\[
\phi(r)=\phi(s).
\]
Then
\begin{align*}
    \phi(r)-\phi(s) &= 0,\\
    \phi(r-s) &= 0.
\end{align*}
So
\[
r-s\in \ker\phi.
\]
Hence
\[
r-s=0,
\]
therefore
\[
r=s.
\]
Thus $\phi$ is injective.
\end{proposition}

\begin{theorem}[Proof of First Isomorphism Theorem]{
Suppose $\phi:\ri\to\si$ is a homomorphism.
Let $K=\ker\phi$, and define
\[
\overline{\phi}(r+K)=\phi(r).
\]
Then $\overline{\phi}$ is a well-defined isomorphism from $\ri/K$ onto $\Ima\phi$.
}
To show $\overline{\phi}$ is well-defined, suppose
\[
r+K=s+K.
\]
Then
\[
r-s\in K=\ker\phi,
\]
so
\[
\phi(r-s)=0.
\]
Hence
\[
\phi(r)=\phi(s).
\]
Thus $\overline{\phi}(r+K)=\overline{\phi}(s+K)$.
Now for $r,s\in \ri$,
\[
\overline{\phi}((r+K)+(s+K))
=\overline{\phi}(r+s+K)
=\phi(r+s)
=\phi(r)+\phi(s)
=\overline{\phi}(r+K)+\overline{\phi}(s+K),
\]
and similarly,
\[
\overline{\phi}((r+K)(s+K))
=\overline{\phi}(rs+K)
=\phi(rs)
=\phi(r)\phi(s)
=\overline{\phi}(r+K)\overline{\phi}(s+K).
\]
So $\overline{\phi}$ is a homomorphism.
It is surjective by definition of $\Ima\phi$.
Indeed, if $s\in \Ima\phi$, then there exists $r\in \ri$ such that $\phi(r)=s$.
Hence
\[
\overline{\phi}(r+K)=s.
\]
Its kernel is trivial.
If
\[
\overline{\phi}(r+K)=0,
\]
then
\[
\phi(r)=0,
\]
so $r\in K$.
Thus
\[
r+K=K=0+K.
\]
Therefore $\overline{\phi}$ is injective.
Hence $\overline{\phi}$ is an isomorphism.
\end{theorem}

\begin{example}{
Prove
\[
\Q[x]/(x^2-2)\cong \Q\left(\sqrt{2}\right).
\]
}
\end{example}

\begin{proof}
Define
\[
\phi:\Q[x]\to \c
\]
by
\[
\phi(f(x))=f\left(\sqrt{2}\right).
\]
Then
\[
\ker\phi \coloneqq \{f(x)\in \Q[x]:f\left(\sqrt{2}\right)=0\}.
\]
Clearly,
\[
x^2-2\in \ker\phi.
\]

\textbf{Claim.}
\[
\ker\phi=(x^2-2).
\]
Indeed, if $f\left(\sqrt{2}\right)=0$, then by the factor theorem over $\Q\left(\sqrt{2}\right)$, $(x-\sqrt{2})$ divides $f(x)$.
Since the coefficients of $f(x)$ are rational, $(x+\sqrt{2})$ also divides the conjugate relation, so
\[
\left(x-\sqrt{2}\right)\left(x+\sqrt{2}\right)=x^2-2
\]
divides $f(x)$ in $\Q[x]$.
Thus every element of $\ker\phi$ lies in $(x^2-2)$, and the reverse inclusion is clear.
Also,
\[
\Ima\phi=\Q\left(\sqrt{2}\right).
\]
Therefore, by the first isomorphism theorem,
\[
\Q[x]/(x^2-2)\cong \Q\left(\sqrt{2}\right).
\]
\end{proof}
